\section*{Capítulo \ref{ch:g11:magnetic-fields} --- El magnetismo y el campo magnético}

\begin{solution}{exo:g11:magnetic-fields:1}
Cada trozo es un \omterm{def:g11:magnetic-fields:magnet}{imán} completo, con un polo norte y otro sur: en cada
cara cortada aparece un polo nuevo. Las caras recién cortadas son un N y
un S, así que se atraen --- los trozos intentan recomponerse. Ninguna
forma de cortarlo aísla un polo: cada corte crea una pareja.
\end{solution}

\begin{solution}{exo:g11:magnetic-fields:2}
(a) Falso: el \omterm{def:g11:electric-gravitational-fields:field}{campo} tiene una sola dirección en cada punto. (b)
Verdadero. (c) Verdadero: cada línea se cierra por el cuerpo del \omterm{def:g11:magnetic-fields:magnet}{imán}.
(d) Falso: la aguja se asienta \emph{a lo largo} de la línea, tangente a
ella.
\end{solution}

\begin{solution}{exo:g11:magnetic-fields:3}
$B = \num{2e-7} \times 10 / 0.020 = \qty{1.0e-4}{T}$ --- el doble del
\omterm{def:g11:electric-gravitational-fields:field}{campo} terrestre.
\end{solution}

\begin{solution}{exo:g11:magnetic-fields:4}
$n = 800/0.40 = \qty{2000}{m^{-1}}$;
$B = \num{1.26e-6} \times 2000 \times 1.5 \approx \qty{3.8}{mT}$. Como
$B \propto I$, duplicar la intensidad da \qty{7.5}{mT}.
\end{solution}

\begin{solution}{exo:g11:magnetic-fields:5}
$F = ILB = 5.0 \times 0.10 \times 0.20 = \qty{0.10}{N}$ --- el \omterm{def:g10:universal-gravitation:weight}{peso} de
$m = F/g = 0.10/9.81 \approx \qty{10}{g}$.
\end{solution}

\begin{solution}{exo:g11:magnetic-fields:6}
\emph{1.} Su derrota apunta $\ang{8}$ al oeste del norte verdadero:
$2000 \times \sin\ang{8} \approx \qty{2.8e2}{m}$ hacia el oeste.

\emph{2.} Gobernar $\ang{8}$ al \emph{este} del norte de la aguja.
\end{solution}

\begin{solution}{exo:g11:magnetic-fields:7}
$d = \mu_0 I / (2\pi B) = \num{2e-7} \times 20 / \num{5.0e-5}
= \qty{8.0}{cm}$. A menos de un decímetro de un cable así, la brújula
lee el hilo y no el planeta --- hay que mantenerla lejos de conductores
con corriente.
\end{solution}

\begin{solution}{exo:g11:magnetic-fields:8}
$n = B/(\mu_0 I) = 0.010 / (\num{1.26e-6} \times 4.0)
\approx \qty{2.0e3}{m^{-1}}$; en \qty{25}{cm},
$N = 1989 \times 0.25 \approx 500$ espiras.
\end{solution}

\begin{solution}{exo:g11:magnetic-fields:9}
$F = 10 \times 0.12 \times 0.50 = \qty{0.60}{N}$;
$a = F/m = 0.60/0.020 = \qty{30}{m/s^2}$ --- unas $3g$: la varilla salta
de verdad.
\end{solution}

\begin{solution}{exo:g11:magnetic-fields:10}
Equilibrio: $ILB = mg$, luego
$I = \num{5.0e-3} \times 9.81 / (0.20 \times 0.10) \approx
\qty{2.5}{A}$. La \omterm{def:g10:inertia:force}{fuerza} tiene que apuntar hacia arriba; con $\vect B$
horizontal y perpendicular al hilo, la mano derecha extendida (con la
palma hacia arriba) fija el sentido necesario de la corriente.
\end{solution}

\begin{solution}{exo:g11:magnetic-fields:11}
Mano derecha extendida: (a) hacia la parte superior de la página; (b)
hacia la derecha; (c) ninguna \omterm{def:g10:inertia:force}{fuerza} --- el hilo es paralelo al \omterm{def:g11:electric-gravitational-fields:field}{campo}.
\end{solution}

\begin{solution}{exo:g11:magnetic-fields:12}
\emph{1.} $B = \num{2e-7} \times 5.0 / 0.030 \approx
\qty{33}{\micro T}$; las líneas rodean el hilo, así que justo debajo de
él el \omterm{def:g11:electric-gravitational-fields:field}{campo} es horizontal y perpendicular al hilo: este--oeste.

\emph{2.} La aguja sigue la suma vectorial:
$\tan\theta = 33/20 = 1.67$, luego $\theta \approx \ang{59}$ respecto
del norte.
\end{solution}

\begin{solution}{exo:g11:magnetic-fields:13}
\emph{1.} $I = U/R = 12/2.0 = \qty{6.0}{A}$;
$n = 400/0.20 = \qty{2000}{m^{-1}}$;
$B_0 = \num{1.26e-6} \times 2000 \times 6.0 \approx \qty{15}{mT}$; con
el \omterm{def:g11:fundamental-interactions:ladder}{núcleo}, $100 B_0 \approx \qty{1.5}{T}$.

\emph{2.} $P = UI = 12 \times 6.0 = \qty{72}{W}$.

\emph{3.} Sin corriente no hay \omterm{def:g11:electric-gravitational-fields:field}{campo} --- y el hierro dulce no se queda
imanado: la carga se suelta al instante, tal como se pretende.
\end{solution}

\begin{solution}{exo:g11:magnetic-fields:14}
\emph{1.} $F = NILB = 50 \times 2.0 \times 0.050 \times 0.30 =
\qty{1.5}{N}$ en cada lado. La corriente recorre los dos lados en
sentidos opuestos, así que las \omterm{def:g10:inertia:force}{fuerzas} son opuestas: un par que hace
girar la espira.

\emph{2.} Cuando es paralelo a $\vect B$, un hilo no nota ninguna \omterm{prop:g11:magnetic-fields:laplace}{fuerza
de Laplace}.

\emph{3.} Media vuelta después, el par se invertiría y deshacería el
giro; el \omterm{rem:g11:magnetic-fields:motor}{conmutador} cambia el sentido de la corriente cada media vuelta
para que la espira siga girando en el mismo sentido.
\end{solution}

\begin{solution}{exo:g11:magnetic-fields:15}
\emph{1.} $d = \num{2e-7} \times 100 / 1 = \qty{2.0e-5}{m} =
\qty{20}{\micro m}$ --- muy \emph{dentro} del cobre, que tiene
milímetros de grosor. Ningún punto accesible cerca de un solo hilo llega
a \qty{1}{T}: los \omterm{def:g11:electric-gravitational-fields:field}{campos} intensos apilan miles de espiras (un
\omterm{def:g11:magnetic-fields:solenoid}{solenoide}), más hierro o superconductores.

\emph{2.} $\num{e8} / \num{5e-5} = \num{2e12}$: unas doce \omterm{def:g10:energy-conservation:power}{potencias} de
diez.
\end{solution}

\begin{solution}{pb:g11:magnetic-fields:1}
\textbf{1.} La aguja es ella misma un \omterm{def:g11:magnetic-fields:magnet}{imán} pequeño; el \omterm{def:g11:electric-gravitational-fields:field}{campo} terrestre
la alinea, con el extremo norte a lo largo de la componente horizontal
de $\vect B$, es decir, hacia el norte magnético.

\textbf{2.} $B = \sqrt{20^2 + 44^2} = \sqrt{2336} \approx
\qty{48}{\micro T}$, inclinado $\tan^{-1}(44/20) \approx \ang{66}$ por
debajo de la horizontal.

\textbf{3.} La derrota va $\ang{6}$ al oeste del norte verdadero:
$30 \times \sin\ang{6} \approx \qty{3.1}{km}$ al oeste.

\textbf{4.} Los errores se suman: $\ang{10}$, luego
$30 \times \sin\ang{10} \approx \qty{5.2}{km}$ --- cinco kilómetros por
diez grados.

\textbf{5.} La componente horizontal: cerca del polo el \omterm{def:g11:electric-gravitational-fields:field}{campo} es casi
vertical, la parte horizontal se encoge hacia cero y a la aguja no le
queda nada con lo que alinearse.

\textbf{6.} $n = 600/0.30 = \qty{2000}{m^{-1}}$.

\textbf{7.} $I = U/R = 12/1.5 = \qty{8.0}{A}$.

\textbf{8.} $B_0 = \mu_0 n I = \num{1.26e-6} \times 2000 \times 8.0
\approx \qty{20}{mT}$.

\textbf{9.} $50 \times \qty{20}{mT} = \qty{1.0}{T}$ --- el doble del
\omterm{def:g11:electric-gravitational-fields:field}{campo} en el polo de un \omterm{def:g11:magnetic-fields:magnet}{imán} de neodimio, y sobre una cara mucho mayor.

\textbf{10.} $P = UI = 12 \times 8.0 = \qty{96}{W}$; en
\qty{10}{\minute}: $E = 96 \times 600 \approx \qty{58}{kJ}$.

\textbf{11.} Porque suelta: se abre el interruptor, el \omterm{def:g11:electric-gravitational-fields:field}{campo} muere y el
coche cae exactamente donde se quiere. Un \omterm{def:g11:magnetic-fields:magnet}{imán} permanente agarra para
siempre.

\textbf{12.} $F = NILB = 100 \times 3.0 \times 0.040 \times 0.25 =
\qty{3.0}{N}$.

\textbf{13.} Un par: aplicadas a ambos lados del eje, las dos \omterm{def:g10:inertia:force}{fuerzas}
opuestas hacen girar la espira en el mismo sentido. Dos \omterm{def:g10:inertia:force}{fuerzas} iguales
y \emph{paralelas} solo empujarían la espira de lado, sin hacerla girar.

\textbf{14.} Ninguna: un hilo paralelo al \omterm{def:g11:electric-gravitational-fields:field}{campo} no nota \omterm{prop:g11:magnetic-fields:laplace}{fuerza de
Laplace}.

\textbf{15.} Invierte la corriente de la espira en cada media vuelta,
justo cuando el par cambiaría de signo --- de modo que el momento
impulsa siempre el mismo giro.

\textbf{16.} Duplicar $N$, o duplicar $I$, o duplicar $B$: la pareja de
\omterm{def:g10:inertia:force}{fuerzas} $F = NILB$ es proporcional a cada uno.

\textbf{17.} $1.5 / \num{5.0e-5} = \num{3.0e4}$: treinta mil veces la
Tierra.

\textbf{18.} $n = B/(\mu_0 I) = 1.5 / (\num{1.26e-6} \times 500)
\approx \qty{2.4e3}{m^{-1}}$ --- y cada una de esas espiras lleva
\qty{500}{A}.

\textbf{19.} $P = RI^2 = 0.20 \times 500^2 = \qty{50}{kW}$ --- la
calefacción de un barrio. El bobinado real es superconductor:
\omterm{def:g11:circuits-and-power:resistance}{resistencia} nula, disipación nula, siempre que se mantenga a unos pocos
grados del cero absoluto.

\textbf{20.} Brújula $\num{5e-5}$, grúa $1$, resonancia $1.5$: el barco
está cinco peldaños por debajo de las dos máquinas, que comparten una
misma década --- entre el $\num{e-10}$ interestelar y el $\num{e8}$ de
la estrella de neutrones. No ha cambiado nada más que los \omterm{def:g11:circuits-and-power:current}{amperios} (y
las espiras, el hierro y el superconductor que los multiplican): un solo
\omterm{def:g11:electric-gravitational-fields:field}{campo}, tres trabajos.
\end{solution}
